\documentclass[11pt,a4paper]{article}
\usepackage[T1]{fontenc}
\usepackage{lmodern,amsmath,amssymb,amsthm,mathtools,booktabs,microtype}
\usepackage[margin=27mm]{geometry}
\usepackage[colorlinks=true,linkcolor=blue,citecolor=blue,urlcolor=blue]{hyperref}
\newtheorem{theorem}{Theorem}[section]
\newtheorem{lemma}[theorem]{Lemma}
\newtheorem{proposition}[theorem]{Proposition}
\newtheorem{corollary}[theorem]{Corollary}
\newcommand{\Q}{\mathbb Q}\newcommand{\Z}{\mathbb Z}
\DeclareMathOperator{\Tr}{Tr}\DeclareMathOperator{\rad}{rad}
\setlength{\emergencystretch}{2em}
\title{Thirteen Variables from Direct Elliptic Index Comparison}
\author{}\date{10 November 2026}
\begin{document}\maketitle
\begin{abstract}
Under the weak rational $abc$ inequality with exponent $121/100$
and the stated elliptic, pole-realization, and prime-pattern inputs,
we construct a computable family of integer polynomials in thirteen
variables, with degree at most $11774654098433$.
The elliptic comparison recovers the input itself as an integer
index, while a separate positive quadratic expression controls its
height. This permits a sixth-power comparison scale. The same
construction with a fourth-power scale has degree at most
$11212784195585$ under the weak inequality with exponent $16/15$.
Ordinary Szpiro implies the first hypothesis; ordinary rational
$abc$ implies both. Every sufficiently large fixed odd coefficient
shift works, but the sufficient thresholds are ineffective.
The input zero is included by a single polynomial factor, without
an additional variable. Complete finite coefficient specifications,
independent exact conversion checks, and uniform degree and support
bounds accompany the proofs.
\end{abstract}

\section{The families and their hypotheses}

All variable counts include the input. For a reduced rational number
$u/v$, with $v>0$, write
\[
 h(u/v)=\log\max(|u|,v).
\]
Fix one of the two pairs
\[
             (\kappa,d)=(121/100,6)\quad\hbox{or}\quad(16/15,4).
\]
The additional height hypothesis is that a constant $K_\kappa\ge1$
exists such that
\begin{equation}\label{eq:weak-abc}
 c\le K_\kappa\operatorname{rad}(abc)^\kappa
 \qquad(a+b=c,\ a,b,c>0\text{ pairwise coprime}).
\end{equation}
Moving the negative terms to the other side gives the corresponding
statement for a coprime nonzero signed triple of sum zero, with the
largest absolute value on the left.

Set
\[
 K=\Q(\sqrt2),\quad D=(3-\sqrt2)^2,\quad
 E:y^2=x^3-Dx,
\]
\[
 P=\left(\frac{57}{8}-\frac72\sqrt2,
                 -\frac{165}{8}+\frac{417}{32}\sqrt2\right).
\]
In addition to \eqref{eq:weak-abc}, retain the following arithmetic inputs with
exactly their stated scope in~\cite{Dependencies,Source27}.
\begin{enumerate}
\item $2E(K)=\Z P$, with the indicated non-torsion generator
and the source index-height bound. The direct point chart below
therefore describes the nonzero multiples of $P$.
\item The pole system, for $b\ne0$ and $B_0=17b$,
\[
 p^2=B_0(1-e)(1+3e),\qquad
 B_0q^2=(s^2-B_0e(B_0-b^2))(s^2+3B_0e(B_0-b^2)),
\]
excludes negative odd valuations of $b$ outside $\{2,3,17\}$.
For a prepared input
\begin{equation}\label{eq:prepared}
 \begin{gathered}
 0<b<17,\quad b\text{ has square reduced denominator},\quad
 b\in\Q_2^{\times2}\cap\Q_3^{\times2}\cap\Q_{17}^{\times2},\\
 v_2(b)\ge3,\quad v_3(b)\ge1,\quad v_{17}(b)\ge2,
 \end{gathered}
\end{equation}
the retained restricted-converse and descent hypotheses supply a
witness with $pqs\ne0$. The nonzero-witness consequence and its
single conic chart are proved in~\cite[Lemma~2]{Source27}.
\item The simultaneous prime-pattern input supplies the specified
prime contact numerators, Gaussian congruences, conjugate reduction
orders, fixed divisibility and avoidance, and compatible real
approximation for a fixed finite list of rational contact coefficients.
Its precise quantified version and the rational-contact and fixed-avoidance
consequences of Sun's theorem are given in~\cite{Dependencies}.
The even number of contacts below permits arbitrary positive rational
scales: Section~\ref{sec:complete} checks the square factors directly.
The order of choices is respected there.
\end{enumerate}
No modular-realization or polynomial contact-height premise is used.
Lemma~\ref{lem:contact} supplies the contact-height estimate from
\eqref{eq:weak-abc}.

The proved one-point case of the aggregate trace-height theorem
\cite{Trace} gives constants $A>0$, $J\ge0$, depending only on
$E,P$, such that
\begin{equation}\label{eq:traceheight}
 n^2\le A h(T(n))+J,\qquad
 T(n)=\Tr_{K/\Q}x([n]P)\quad(n\in\Z\setminus\{0\}).
\end{equation}
This applies von Buhren's theorem after checking the absence of
constant abelian translates for the fixed trace function. Its additive
constant is ineffective. It is a proved height theorem, not an
additional conjecture. For a nonzero multiple of $P$, the abscissa
is a positive square in both real embeddings; the curve equation
then gives $x>3\mp\sqrt2$ in the respective embeddings, hence
$T(n)>6$.

\begin{theorem}\label{thm:main}
Retain the arithmetic inputs above. For each of the two pairs
$(\kappa,d)$, there is a computable family
\[
 F_B^{(d)}\in\Z[X,Y_1,\ldots,Y_{12}],\qquad B\ge3\text{ odd},
\]
with the following properties.
\begin{enumerate}
\item Each member has a complete finite integer coefficient
specification. Uniformly in $B$, its degree and number of ordinary
monomials satisfy
\[
\begin{array}{c|r|c}
 d&\text{degree upper bound}&\text{support upper bound}\\\hline
 6&11774654098433&<2^{508}\\
 4&11212784195585&<2^{507}.
\end{array}
\]
In particular, both degree bounds are less than $2^{44}$.
\item Under \eqref{eq:weak-abc} for the indicated $\kappa$, there
is a constant $B_*^{(d)}$ such that every odd $B\ge B_*^{(d)}$
has rational input projection exactly $\Z$.
\item No numerical sufficient value of $B_*^{(d)}$ is determined.
The supplied $B=3$ members are coefficient-verification fixtures,
not certified defining members.
\end{enumerate}
The fixed coefficient $B$ is never a quantified variable.
\end{theorem}

Ordinary rational $abc$ implies both displayed weak inequalities.
Oesterl\'e proves that ordinary discriminant--conductor Szpiro implies
weak $abc$ with every exponent $6/5+\epsilon$
\cite[Section~I.3, p.~169]{Oesterle}. Taking $\epsilon=1/100$
therefore supplies the sixth-power family, even from the semistable
case of ordinary Szpiro. No implication from ordinary Szpiro to the
stronger exponent $16/15$ is asserted. None of these implications
removes the retained elliptic and pole-realization premises.

\section{Height estimates}

Put
\[
 \Theta=\left\{0,-\frac12,-\frac14,-\frac25,-1,-\frac13\right\},
 \qquad S_0=\{p:p\le76176\}.
\]
For $t\in\Q\setminus\Theta$, let $M_0(t)$ be the product of
primes outside $S_0$ for which $v_p(t-\theta)$ is positive
and odd for some $\theta\in\Theta$. Define
\[
                  \gamma=8/\kappa-6,\qquad \alpha=d/\kappa-3.
\]
Both numbers are positive for either pair used here.

\begin{lemma}\label{lem:contact}
Under \eqref{eq:weak-abc}, there is a constant $C_0\ge0$, independent
of every later coefficient shift, such that
\[
                 h(t)\le\gamma^{-1}\log M_0(t)+C_0
                  \qquad(t\in\Q\setminus\Theta).
\]
\end{lemma}
\begin{proof}
Consider the degree-four rational function
\[
 \varphi(z)=\frac{z(z+2)^3}{(2z+1)^3},\qquad
 \varphi(z)-1=\frac{(z-1)^3(z+1)}{(2z+1)^3}.
\]
Its complete inverse image of $\{0,1,\infty\}$ is
$\{0,-2,-1/2,1,-1,\infty\}$. For a primitive integer pair
$(x,y)$ away from these six points, put
\[
 A=x(x+2y)^3,\quad C=y(2x+y)^3,\quad
 A-C=(x-y)^3(x+y),
\]
\[
                G(x,y)=xy(x+2y)(2x+y)(x-y)(x+y).
\]
The polynomial identity for $A-C$ follows by expansion.
A prime dividing both $A$ and $C$ must divide both $x+2y$
and $2x+y$, so it is $3$. At that prime $x,y$ are units
and the smaller of those two valuations is one, because
$2(x+2y)-(2x+y)=3y$. Hence $\gcd(A,C)$ is either $1$ or $27$.
After division by this gcd, the three integers $A,C,A-C$
are pairwise coprime, and their radical is supported on $G$.

Write $H=\max(|x|,|y|)$. One has
$\max(|x+2y|,|2x+y|)\ge H$. If the corresponding prefactor
$|x|$ or $|y|$ is at least $H/4$, this gives
$\max(|A|,|C|)\ge H^4/4$. Otherwise the other prefactor
is $H$ and its linear factor is at least $H/2$, giving
$\max(|A|,|C|)\ge H^4/8$. Consequently
\[
                    h(\varphi(x/y))\ge4\log H-\log216.
\]
The weak inequality therefore implies
\[
 \log\operatorname{rad}(G(x,y))
       \ge\frac4\kappa\log H-\frac{\log216+\log K_\kappa}{\kappa}.
\]

For a nonzero integer $N$, let $\operatorname{sf}(N)$ be the
product of primes of odd valuation in $N$. Prime by prime,
$\operatorname{rad}(N)^2\le |N|\operatorname{sf}(N)$.
Also $|G(x,y)|\le36H^6$. Outside the fixed set $S_0$, two
of its six primitive linear factors cannot vanish together.
Thus its squarefree part outside $S_0$ is precisely the odd
contact radical for the six displayed points. The contribution
from $S_0$ is bounded by the fixed product of its primes.
Combining these estimates gives
\[
 \log M(z)\ge\left(\frac8\kappa-6\right)h(z)-C
             =\gamma h(z)-C.
\]
Finally $z=2+1/t$ sends the six contacts in $\Theta$ to the
six points above. Its integral matrix has determinant $-1$;
its inverse is integral too, and both matrices have maximum
row sum three. Hence it changes logarithmic height by at most
$\log3$. Substitution of a primitive numerator-denominator pair
shows that the odd contact radicals agree outside $S_0$.
This proves the lemma after adjusting the fixed constant.
\end{proof}

For $u,v\in\Q$, let
\[
 h_2(u,v)=h(1:u:v)=\log\max(D_0,|A'|,|B'|),
\]
where $D_0>0$, $u=A'/D_0$, $v=B'/D_0$, and
$\gcd(D_0,A',B')=1$. Each individual height is at most $h_2$.

\begin{lemma}\label{lem:powers}
For an even $d$ with $\alpha=d/\kappa-3>0$, one has
\[
                  h_2(u,v)\le\alpha^{-1}h(u^d+v^d)+C_1
                     \qquad(u,v\in\Q),
\]
where one may take
\[
 C_1=\frac{\alpha+1}{d\alpha}\log2
                  +\frac1{\kappa\alpha}\log K_\kappa.
\]
\end{lemma}
\begin{proof}
If one coordinate is zero, $h_2(u,v)=h(u^d+v^d)/d$;
if both are zero all heights vanish. Otherwise write
$g=\gcd(|A'|,|B'|)$, $A_0=A'/g$, and $B_0=B'/g$.
Then $\gcd(A_0,B_0)=\gcd(g,D_0)=1$. Put
\[
 H=\max(|A_0|,|B_0|),\qquad N=A_0^d+B_0^d,
 \qquad H^d\le N\le2H^d.
\]
Set $k=\gcd(N,D_0^d)$, $m_0=N/k$, and $n_0=D_0^d/k$.
The reduced fraction for $u^d+v^d$ is $m/n_0$, where
$m=g^d m_0$. Let $M=\max(m,n_0)$, so its height is $\log M$.
Every prime dividing $N$ either divides $D_0$ or survives in
$m_0$, hence $\operatorname{rad}(N)\le D_0m_0$.
Applying \eqref{eq:weak-abc} to the coprime positive triple
$A_0^d,B_0^d,N$ gives
\[
               H^{d/\kappa-2}\le K_\kappa^{1/\kappa}D_0m_0.
\]
Since $D_0^d=n_0N/m_0\le2n_0H^d/m_0$, this yields
\[
 H^\alpha\le2^{1/d}K_\kappa^{1/\kappa}
                  m_0^{(d-1)/d}n_0^{1/d}.
\]
Also $g=(m/m_0)^{1/d}$, so
\[
                \max(gH,D_0)\le2^{1/d}H(M/m_0)^{1/d}.
\]
Combining the logarithms proves
\begin{align*}
 h_2(u,v)
 &\le C_1+\frac{d-1-\alpha}{d\alpha}\log m_0
          +\frac1{d\alpha}\log n_0+\frac1d\log M\\
 &\le C_1+\alpha^{-1}\log M.
\end{align*}
The three coefficients are nonnegative and sum to $1/\alpha$;
indeed $d-1-\alpha=d(1-1/\kappa)+2>0$.
This calculation includes the common numerator factor and all
cancellation with the denominator.
\end{proof}

\begin{lemma}\label{lem:input-height}
For every rational $X$ and integer $B\ge3$, put
\begin{equation}\label{eq:input}
                              \beta=X^2+B.
\end{equation}
Then
\[
                       \beta\ge B,\qquad 2h(X)\le h(\beta).
\]
\end{lemma}
\begin{proof}
Write $X=u/v$ in lowest terms with $v>0$. Since
$\gcd(u^2+Bv^2,v^2)=1$, the displayed numerator and denominator
of $\beta=(u^2+Bv^2)/v^2$ are already reduced. Moreover
\[
                     u^2+Bv^2\ge\max(u^2,v^2).
\]
Since $\beta\ge B>1$, its height is the logarithm of this
positive numerator. This proves the height inequality, including
$u=0$. There is no cancellation factor.
\end{proof}

The integer label compared with the elliptic index is $X$.
The separate quantity $\beta$ appears only in the height scale;
there is no need to infer integrality of $X$ from integrality
of $\beta$.

\section{The complete rational source}

Fix an odd integer $B\ge3$ and use \eqref{eq:input}.
Retain three pairs of rational coordinates, one for each of the
points $A,H,U$. For any such pair $(h_0,h_1)$ put
\[
 h=h_0+h_1\sqrt2,\quad x=h^2,\quad j^2=h^4-D,\quad y=hj.
\]
The direct point-chart theorem of~\cite{Source27} proves that
these are exactly the nonzero multiples of $P$. Each finite
choice $j$ is encoded by one scalar monic quartic as follows.
Write $h^4-D=a+b\sqrt2$, put
$c=2a+b$, $d=-a-b$, and write
$(c+d\sqrt2)^2(a+b\sqrt2)=a'+b'\sqrt2$.
Then
\[
 Y^4-4a'Y^2+8(b')^2=0,\quad
 j=\frac{Y/2+(b'/Y)\sqrt2}{c+d\sqrt2}.
\]
Its discriminant is
$2^{15}(b')^2(2a^2-b^2)^4(a^2-2b^2)^2$, nonzero at every
rational pair $(h_0,h_1)$. Indeed $h^4-D\ne0$ by its norm,
the two quadratic norm factors are anisotropic over $\Q$, and
$b'=-4a^3+6ab^2+3b^3$ is anisotropic by Eisenstein at 3.
The square-root chart is complete because its transformed rational
component is $2(u^3-2v^3)$ for $j=u+v\sqrt2$, which is nonzero
for $j\ne0$. These are algebraic facts independent of the height hypothesis.

Compute $V_0=H-P$ and $V_1=U-A$ by the usual affine addition
formulas. Put $f=1/x$, $z=-x/y$, let $\sigma$ be conjugation
of $K/\Q$, and define
\[
 L=\frac{z(U)}{z(H)},\qquad
 L'=\sigma\!\left(\frac{z(V_1)}{z(V_0)}\right).
\]
Use precisely the six field-valued tests
\begin{equation}\label{eq:tests}
 D_1=f(H),\ D_2=f(U),\ D_3=\sigma f(V_0),\
 D_4=\sigma f(V_1),\ D_5=L-L',\ D_6=L-X.
\end{equation}
Every test is assigned order one. No scalar-height pole or Lucas
relation is present.

Let $T=\Tr x(A)$, and use the two positive rational scales
\[
                         a_0=1,\qquad a_1=\beta^d+T^d.
\]
Retain one rational slope $\lambda$. For $r=0,1$, set
\begin{align}
 t_r&=a_r\lambda,&\delta_r&=t_r-1,&
 F_r&=\prod_{\theta\in\Theta}(t_r-\theta),&
 R_r&=\frac{F_r}{\delta_r^8},\label{eq:detector}\\
 W_r&=\prod_{j=1}^6\left(1+
             \frac{D_j^2}{\delta_r^4R_r}\right),&
 b_r&=\frac{R_r}{\Tr(W_r^{-1})^2}.\label{eq:W}
\end{align}
All displayed divisions, affine-chart denominators, $F_r$,
$\delta_r$, $\operatorname{Norm}(W_r)$, and $\Tr(W_r^{-1})$
are guarded nonzero. These conditions are part of the source.

For each $b=b_r$, retain independent rational parameters $\tau_r,s_r$.
Put $B_0=17b$ and
\[
 e=\frac{\tau_r^2-B_0}{\tau_r^2+3B_0},\qquad
 q^2=\frac{(s_r^2-B_0e(B_0-b^2))
                  (s_r^2+3B_0e(B_0-b^2))}{B_0}.
\]
Guard the denominator and require $q\ne0$. This reconstructs
the first pole ordinate as
$p=-4B_0\tau_r/(\tau_r^2+3B_0)$.
Each rational-square equation $q^2=N_0/D_0$ is made monic
by $Z=D_0q$, giving $Z^2=N_0D_0$ with $Z\ne0$ and $D_0\ne0$.
No new unrestricted rational parameter is used for these roots.

\section{The local implication}

\begin{lemma}\label{lem:local}
For any rational solution of the guarded source, write
$A=[n]P$, $H=[h]P$, $U=[u]P$.
If $X\ne n$ and $N$ is the positive absolute numerator
of $X-n$ in lowest terms, then
\[
                         M_0(t_r)\mid N\qquad(r=0,1).
\]
\end{lemma}
\begin{proof}
Fix a prime $p\mid M_0(t_r)$ and abbreviate $R=R_r$, $W=W_r$.
Outside $S_0$ only one contact factor has positive valuation,
$\delta_r$ is a unit, and $v_p(R)$ is odd because its denominator is an eighth power. For each place $\mathfrak p$ above $p$, the two terms
of a factor of $W$ have valuations of different parity. Hence
\[
 v_{\mathfrak p}(W)=\sum_j
       \min\{0,2v_{\mathfrak p}(D_j)-v_p(R)\}\le0.
\]
Put $w=\max_{\mathfrak p\mid p}v_{\mathfrak p}(W)\le0$.
Then $v_p\Tr(W^{-1})\ge-w$.
The valuation of $b=R/\Tr(W^{-1})^2$ is odd, so the pole
obstruction gives
\[
                   0<v_p(b)\le v_p(R)+2w.
\]
Thus $\nu=v_p(R)>0$ and $w>-\nu/2$. At a place attaining
$w$, each summand in the valuation of $W$ is at least $w$.
It follows that $v_{\mathfrak p}(D_j)>\nu/4>0$, so all
six tests have valuation at least one at that place.

The first four tests put $H,U$ into the formal group at
$\mathfrak p$ and $V_0,V_1$ into it at the conjugate place.
At these good unramified places the formal logarithm satisfies
\[
                         \log_E(Q)=z(Q)(1+O(p^4)).
\]
Indeed its first omitted term is a fixed multiple of $Df(Q)^2$;
$v_p(f(Q))\ge2$ in the formal group. The remaining coefficient
estimates, uniformly for $p>76176$, are those proved in the
formal-logarithm paper~\cite{Logs}. The exact logarithmic ratios
are $u/h$ and $(u-n)/(h-1)$, respectively.

The affine guards imply $h\ne0,1$. At least one of $h,h-1$
is a $p$-adic unit, so the corresponding exact ratio is integral.
Relative precision four makes its truncated ratio integral.
The test $L-L'\equiv0\pmod p$ makes the other truncated ratio
integral as well; the same relative precision then makes the
other exact ratio integral. Relative errors are now absolute
errors of valuation at least four. Consequently
\[
 L\equiv u/h,\qquad L'\equiv(u-n)/(h-1),\qquad
 L\equiv L'\pmod p.
\]
Multiplying by the integral numbers $h,h-1$ and subtracting yields
$L\equiv n\pmod p$. The target test gives $L\equiv X\pmod p$.
In particular $X$ is $p$-integral, and $p$ divides the
reduced numerator of $X-n$. This holds at every consumed
prime and proves the divisibility. The same argument includes
inert primes; if their formal-group requirements are incompatible,
no source solution occurs at that prime.
\end{proof}

\section{Soundness and a uniform sufficient shift}

Suppose a rational source solution has $X\ne n$, and let $N$
be the positive absolute numerator of $X-n$. The point chart gives
$n\ne0$, because $A$ is a nonzero multiple of the non-torsion point $P$.
Lemmas~\ref{lem:local} and~\ref{lem:contact} give
\[
 h(\lambda),h(a_1\lambda)\le\gamma^{-1}\log N+C_0.
\]
The nonzero slope and the quotient-height inequality therefore imply
\[
                       h(a_1)\le2\gamma^{-1}\log N+2C_0.
\]
Apply Lemma~\ref{lem:powers} to $\beta,T(n)$ and set
\[
 \eta=h_2(\beta,T(n)),\qquad q=\frac2{\gamma\alpha},\qquad
 C=\frac{2C_0}{\alpha}+C_1.
\]
Then
\begin{equation}\label{eq:combined}
                            \eta\le q\log N+C.
\end{equation}
The constants depend on the fixed contact function, $d$, and the
weak-$abc$ constant, but not on $B$, $X$, or any witness.
The relevant rational numbers are
\[
\begin{array}{c|c|c|c|c}
(\kappa,d)&\gamma&\alpha&q&q/2\\\hline
(121/100,6)&74/121&237/121&14641/8769&14641/17538\\
(16/15,4)&3/2&3/4&16/9&8/9.
\end{array}
\]
In both cases $q/2<1$.

Writing $X=u/v$ in lowest terms gives
\[
 N=|u-nv|\le e^{h(X)}(1+|n|),\qquad
 h(X)\le\frac12h(\beta)\le\frac12\eta.
\]
Consequently \eqref{eq:combined} implies
\[
 \eta\le c\log(1+|n|)+C',\qquad
 c=\frac q{1-q/2},\quad C'=\frac C{1-q/2}.
\]
By the one-point trace-height estimate \eqref{eq:traceheight},
\[
 |n|^2\le Ac\log(1+|n|)+AC'+J
                  \le Ac|n|+AC'+J.
\]
Choose an integer $M>Ac+\sqrt{AC'+J}$. It follows that $|n|<M$,
and hence
\[
                   h(\beta)\le\eta
                  \le c\log(1+M)+C'=:H_*.
\]
If $B_*>e^{H_*}$, every odd $B\ge B_*$ contradicts
\[
                      h(\beta)\ge\log\beta\ge\log B>H_*.
\]
Thus every rational source solution for such a shift has
$X=n\in\Z\setminus\{0\}$. The proof is uniform in the shift;
it supplies no effective numerical value of $B_*$.

\section{Completeness for every nonzero integer}\label{sec:complete}

Fix $X=n\in\Z\setminus\{0\}$ and an odd $B\ge3$.
Fix $A=[n]P$ and the two scales before choosing any prime.
The scales are positive and distinct, since $\beta\ge3$ and
$d$ is even. Temporarily put $H=[h]P$ and $U=[hn]P$.
The sign and parity of $n$ are unrestricted.

Choose a nonexceptional $h_*\in64\Z$, and work in the real
orbit closure of $64\Z P$. Choose a compact real neighborhood
for $H$ around $[h_*]P$ on which the point coordinates, both difference charts,
and the six tests are defined and bounded. Such a neighborhood
exists: the rational functions restricted to the multiplication
curve $H\mapsto(H,[n]H)$ are not identically singular.
For every nonzero integer $n$, the multiplication map $[n]$ is
nonconstant. The forbidden equalities $U=\pm A$ are equivalent
to $h=\pm1$, even when $n$ is negative or even. Similarly,
$H=\pm P$ means $h=\pm1$. Excluding $h=0,\pm1$ therefore
handles both difference charts and all identity points. Every
other nonzero multiple of the non-torsion $P$ has nonzero
abscissa and ordinate.
Pass to a smaller compact neighborhood to avoid all remaining
zero denominators. Enlarge $L$ with $64\mid L$ and the real
component-group condition so that the progression
$h\equiv h_*\pmod L$
preserves fixed lower valuation bounds for these tests at
$2,3,17$. These are precisely the local-continuity and real-orbit
steps of the source realization, with fewer point roles and no
scalar constraint.

Choose a fixed $D_0$ with $D_0^2\theta/a_r\in\Z$ for all
contacts and both scales, set $v=D_0^2w^2$, and seek
$\lambda=u/v$. Apply the retained prime-pattern theorem to
the distinct forms $u-(\theta/a_r)v$, after repetitions have
been removed. All bad primes and previously fixed moduli are
excluded before this application. The theorem permits arbitrarily
large divisibility of $w$, a sufficiently large prescribed lower
bound on $u/w^2$, and the conjugate reduction-order conditions.

As in~\cite{Dependencies,Baseline}, the annihilating order moduli
$U_0,V_0$ for the two conjugate families satisfy
\[
 \gcd(U_0,V_0)=1,\qquad \gcd(U_0,L)=16,\qquad\gcd(V_0,L)=1.
\]
Thus the congruences
$h\equiv0\pmod{U_0}$, $h\equiv1\pmod{V_0}$,
$h\equiv h_*\pmod L$ are compatible, since $16\mid h_*$.
Choose $h$ in that
progression while retaining the real neighborhood and avoiding
the finitely many excluded multiples. At one place above each
contact prime, $H,U$ are formal; at the conjugate place,
$H-P,U-A$ are formal. Both exact logarithmic ratios are $n$.
The constant approximation has precision four, so all six tests
have their assigned order one, including the direct target $n=X$. The integer $n$ is $p$-integral
for every selected prime; oddness is not used. When $n=1$ or $-1$,
some comparison tests vanish identically, which is permitted: their
factors in $W_r$ equal one.

The inherited detector preparation now applies without alteration.
For clarity, its main ingredients are the following. The identity
\[
 R_r=\frac{a_r^6\prod_{\theta\in\Theta}(u-(\theta/a_r)v)}
                  {v^6(a_r\lambda-1)^8}
\]
has a rational-square prefactor. In particular, the possibly nonsquare
scale contributes the square $a_r^6$.
At each chosen contact prime $v_p(R_r)=1$; the successful
conjugate place gives all inverse norm-factor residues equal to
one, and the trace of $W_r^{-1}$ is a unit with residue one or two.
The resulting $b_r$ has square reduced denominator, and the
prescribed congruences give its local squareclasses at $2,3,17$.
The fixed lower bounds on the tests and sufficient divisibility
of $w$ give the depths in~\eqref{eq:prepared}.
On the compact real set, uniformly,
\[
 R_r\sim a_r^{-2}\lambda^{-2},\qquad
 W_r=1+O(\lambda^{-2}),\qquad
 b_r\sim(4a_r^2\lambda^2)^{-1}.
\]
Thus $0<b_r<17$ for the sufficiently large lower bound on
$\lambda$ fixed before prime selection. All detector guards
are nonzero. These calculations are the source comparison-pole
preparation proof~\cite{Baseline}; they depend on the finite
list of tests and their orders, not on a Lucas or scalar test.

Each prepared input has an independent pole certificate with
$pqs\ne0$ by the retained pole-realization hypotheses.
Since $p\ne0$ excludes $e=1$, the substitution
$\tau=p/(e-1)$ gives the single chart used above and its
nonzero denominator. Thus both nonzero quadratic-root conditions
are realized. The direct point charts cover $A,H,U$.
Every nonzero integer input therefore has all required witnesses,
for every fixed odd $B\ge3$. The input zero is included by the
polynomial operation in the next section.

\section{The polynomial and its exact specification}

The retained rational coordinates are
\begin{center}
\begin{tabular}{lr}
\toprule
Coordinates & Number\\\midrule
Input $X$ & 1\\
Three point pairs for $A,H,U$ & 6\\
Comparison slope $\lambda$ & 1\\
Two independent pole pairs $(\tau_r,s_r)$ & 4\\\midrule
Retained before finite-root elimination & 12\\\bottomrule
\end{tabular}
\end{center}
There are three monic point quartics and two monic pole quadratics.
The point algebra has rank $b=4^3=64$, and the full algebra has
rank $n_0=4^3 2^2=256$. Its guard has the form $gq_1q_2$, with
$g$ and both quadratic radicands in the point algebra. There
are no further residual equations.

The separating-element theorem~\cite{Active} replaces all five
finite roots and every nonzero guard with one rational coordinate
$U_0$. Let $Q_B^{(d)}$ be its exact integer polynomial. Define
\begin{equation}\label{eq:final}
                         F_B^{(d)}=XQ_B^{(d)}.
\end{equation}
For $X\ne0$, the rational zero sets of $F_B^{(d)}$ and $Q_B^{(d)}$
are the same. For $X=0$, equation \eqref{eq:final} is satisfied
by every rational choice of the other coordinates. Soundness and
completeness therefore give exactly $\Z$ for every sufficiently
large fixed odd shift. The operation introduces no new coordinate,
so the total remains thirteen.

For completeness, the finite formula for $Q_B^{(d)}$ is as follows.
Let $B_0$ denote the point algebra, and $A_0$ the algebra after
adjoining $q_1,q_2$, with $q_j^2=c_j\in B_0$. Put
\[
 T_0=g^2c_1c_2,\qquad
 \det_{B_0}(I-vM_{T_0})=\sum_{j=0}^{64}a_jv^j,\qquad a_0=1.
\]
Write $\theta_1,\ldots,\theta_5$ for the native generators and
$\delta_i$ for the discriminant of the $i$th point quartic.
The scalar polynomial
\[
 R_0=256\left(1+\sum_{i=1}^{5}\sum_{u,v=1}^{256}
                  (M_{\theta_i}^{A_0})_{uv}^2
       +\sum_{u,v=1}^{64}(M_g^{B_0})_{uv}^2\right)
\]
bounds all required eigenvalues. For each $1\le r\le64$, set
\begin{align*}
 c_r&=a_r,&h_r(Z)&=\sum_{j=0}^r a_jZ^{r-j},&E_r&=c_r-h_r(T_0),\\
 D_r&=c_r^2\prod_{i=1}^3\delta_i^2,&s_r&=3r+36,\\
 K_r&=(2R_0)^{s_r}(1+c_r^2)\prod_{i=1}^3(1+\delta_i^2),&
 C_r&=D_r+4R_0K_r,\\
 \omega_r&=D_rE_r\left(C_r^5+\sum_{i=1}^5D_r^iC_r^{5-i}\theta_i\right),&
 H_r(U_0)&=\det_{A_0}(I-U_0M_{\omega_r}).
\end{align*}
Then
\[
 Q_B^{(d)}=H_{64}(U_0)\prod_{r=1}^{63}
                 \left(H_r(U_0)^2+\sum_{j=r+1}^{64}a_j^2\right).
\]
The guard idempotents isolate the nonzero part, and the squared tail
coefficients select its correct degree. The separation theorem
covers nonreal points, vanishing branches, and nilpotents on the
inactive part. Thus this is a finite integer coefficient algorithm,
not a prescription to select a favorable degree case.

\subsection{Degree, support, and coefficient bounds}

The independently checked native weights and guard weights are
\[
\begin{array}{c|c|r|r}
 d&\text{generator weights}&w_G&w_g\\\hline
 6&(6,6,6,51093,65097)&142557&26367\\
 4&(6,6,6,51093,60429)&136845&25323.
\end{array}
\]
The separating-element weights have the linear upper bound
$W_1r+W_0$, with
\[
\begin{array}{c|r|r}
 d&W_1&W_0\\\hline
 6&9146712&67163859\\
 4&8718120&63704871.
\end{array}
\]
Summing the principal-minor degree bounds over all 64 factors gives
\[
 \deg Q_B^{(d)}\le256\bigl(W_1\,64^2+(W_0+1)127\bigr)=:\mathcal D_d.
\]
The exact values are
\[
             \mathcal D_6=11774654098432,\qquad
             \mathcal D_4=11212784195584.
\]
Equation \eqref{eq:final} increases the degree by at most one,
which proves the bounds in Theorem~\ref{thm:main}.
The shift is a formal coefficient parameter of weight zero, and
no branch or operation depends on its numerical value. The bounds
therefore hold for every fixed odd $B\ge3$.

The degree in $U_0$ is at most $E=256\cdot127=32512$.
Counting all monomials within these two degree bounds gives
\[
 \#\operatorname{supp}(Q_B^{(d)})
 \le \binom{\mathcal D_d+13}{13}
                    -\binom{\mathcal D_d-E+12}{13}.
\]
Indeed monomials with $U_0$ exponent at least $E+1$ are in
bijection with monomials of degree at most $\mathcal D_d-E-1$
after subtracting that exponent. Multiplication by $X$ is an
injective shift of exponent vectors, so it preserves the exact
support size and the complete coefficient multiset. Exact integer
calculation bounds the displayed expression by $2^{508}$ for
$d=6$, and by $2^{507}$ for $d=4$, with strict inequalities.
The determinant polynomials are not expanded, and no exact support
count is asserted.

For the algebraic fixtures $B=3$, the coefficient-norm certificates give
\[
\begin{array}{c|r}
 d&\log_2\|F_3^{(d)}\|_1\text{ upper bound}\\\hline
 6&28598607557817\\
 4&27756172773305.
\end{array}
\]
Both numbers are below $2^{45}$. These are fixture bounds only;
neither $B=3$ member is certified to define the integers.

\subsection{Reconstruction and its scope}

The module \nolinkurl{computation/few_variables/direct_input13_v22}
contains the source builder, the exact indexed specifications,
and independent source and bound checkers. Its \texttt{weak}
and \texttt{abc} directories contain the sixth- and fourth-power
variants, respectively. Each source has 1756 rational-circuit nodes
and 752 integer coefficient nodes. The exact conversion audit
checks the three point quartics and their derivatives, two rational
square clearings, and all 48 atomic guards. The changed scale and
target identities are separately verified symbolically with $B$
treated as a formal coefficient.

An independent implementation also reconstructs the complete
rational formulas in fifteen finite-field samples per variant.
These samples are bounded supplementary evidence. They do not
replace the symbolic conversion, the general local proof, or
completeness. The degree checker imports no producer, checks
dependencies and companion-column inequalities, and sums all
64 degree cases independently. A final exact check verifies
\eqref{eq:final}, its coordinate count, and its degree, support,
and coefficient consequences. None of these computations verifies
weak $abc$ or the retained arithmetic premises.

\section{The range of weak exponents}

\begin{corollary}
Retain the elliptic, pole, and prime-pattern inputs. Suppose weak
rational $abc$ holds for one exponent $1<\kappa<4/3$. Choose an
even integer $d$ such that
\[
 \left(\frac8\kappa-6\right)
 \left(\frac d\kappa-3\right)>1.
\]
Direct comparison with $X$, the scale
$(X^2+B)^d+T(n)^d$, and the final factor $X$ give a computable
thirteen-variable family defining $\Z$ for all sufficiently large
fixed odd shifts.
\end{corollary}
\begin{proof}
Here $\gamma>0$ and $\gamma\alpha>1$. The height lemmas and
local divisibility give $\eta\le2(\gamma\alpha)^{-1}\log N+O(1)$.
Lemma~\ref{lem:input-height} makes the coefficient of $\eta$ on
the right equal to $(\gamma\alpha)^{-1}<1$. The uniform
soundness proof applies. Evenness of $d$ preserves positivity,
and completeness for nonzero integers is unchanged. Such an
even $d$ exists for every stated $\kappa$.
\end{proof}

For $d=6$, the sufficient exponent range is
\[
 1<\kappa<\frac{30-2\sqrt{21}}{17};
\]
for $d=4$, it is
\[
 1<\kappa<\frac{24-4\sqrt2}{17}.
\]
These follow by solving $\gamma\alpha>1$ in the positive range.
The first contains $121/100$ and the second contains $16/15$.
The input-height separation does not extend the six-contact
estimate beyond $\kappa<4/3$: positivity of $\gamma$ is still
required. The numerical degree and support certificates above
concern precisely the two supplied power scales.

\clearpage
\begin{thebibliography}{9}
\bibitem{Oesterle}
J.~Oesterl\'e, Nouvelles approches du ``th\'eor\`eme'' de Fermat,
\emph{Ast\'erisque} \textbf{161--162} (1988), S\'eminaire
Bourbaki, expos\'e no.~694, 165--186; Section~I.3, p.~169.
\url{https://www.numdam.org/article/SB_1987-1988__30__165_0.pdf}.
\bibitem{Source27}
\emph{A defining polynomial in twenty-seven variables:
an existence theorem with a uniform degree bound},
companion paper,
\nolinkurl{few_variables27_nontorsion_20261109_v15.tex},
Lemma~2 and the direct point chart, 9 November 2026.
\bibitem{Dependencies}
\emph{Arithmetic dependencies of the construction with 29 variables},
companion paper,
\nolinkurl{arithmetic_dependencies_20261108_v8.tex}.
It pins the source statements and Sun's prime-pattern theorem.
\bibitem{Baseline}
\emph{An explicit polynomial in 29 variables over the rationals},
companion paper, \nolinkurl{few_variables_20261108_v8.tex},
local implications and simultaneous-completeness sections.
\bibitem{Trace}
\emph{Aggregate trace heights of elliptic multiples}, companion
paper, \nolinkurl{aggregate_trace_height_20261108_v9.tex},
applying J.~von Buhren,
\emph{J. Th\'eor. Nombres Bordeaux} \textbf{29} (2017), 307--320,
Theorem~1.1, \url{https://www.numdam.org/item/10.5802/jtnb.981.pdf}.
\bibitem{Logs}
\emph{Uniform Rational Approximations to an Elliptic Logarithm},
companion paper, \nolinkurl{elliptic_logarithms_20261108_v8.tex}.
\bibitem{Active}
\emph{Separating the Nonzero Part of a Finite Algebra},
companion paper, \nolinkurl{active_finite_algebra_20261109_v15.tex}.
\end{thebibliography}
\end{document}
